% Chinese draft: compile with XeLaTeX (requires ctex and appropriate CJK fonts).
% For English-only submissions use amsart or the journal's official class.
\documentclass[UTF8,11pt,a4paper]{ctexart}
\usepackage{amsmath,amssymb,amsthm,geometry,booktabs,graphicx}
\usepackage[hidelinks]{hyperref}
\geometry{margin=2.7cm}
\newtheorem{theorem}{定理}[section]
\newtheorem{lemma}[theorem]{引理}
\newtheorem{proposition}[theorem]{命题}
\newtheorem{corollary}[theorem]{推论}
\theoremstyle{definition}
\newtheorem{definition}[theorem]{定义}
\newtheorem{example}[theorem]{例}
\newtheorem{conjecture}[theorem]{猜想}
\theoremstyle{remark}
\newtheorem{remark}[theorem]{注}
\numberwithin{equation}{section}
\title{研究题目}
\author{} % Fill only with author-confirmed information.
\date{}
\begin{document}
\maketitle
\begin{abstract}
% State the problem, approach, proved results, and limitations.
\end{abstract}
% Add verified keywords and MSC classification when available.
\section{引言}
% Literature, gap, exact contributions, main theorem and roadmap.
\section{预备知识}
% Definitions, notation, scope, cited lemmas.
\section{主要结果及证明}
% Split the proof into checkable lemmas; use labels and references.
\section{推论与应用}
\section{讨论与开放问题}
% Data/code availability and AI disclosure as appropriate.
% Insert a bibliography with verified entries only.
\appendix
\section{计算证书与复现说明}
% Include completeness arguments, commands and actual results when relevant.
\end{document}
